% Generated draft: requires mathematical and source review.
\documentclass[11pt]{article}
\usepackage[margin=25mm]{geometry}
\usepackage{fontspec}
\setmainfont{Latin Modern Roman}
\usepackage{amsmath,amssymb,mathtools,unicode-math}
\setmathfont{Latin Modern Math}
\usepackage{enumitem,needspace,xurl,hyperref}
\setlength{\parindent}{0pt}
\setlength{\parskip}{5pt}
\setlength{\emergencystretch}{3em}
\newcommand{\N}{\mathbb{N}}
\newcommand{\Z}{\mathbb{Z}}
\newcommand{\Q}{\mathbb{Q}}
\newcommand{\R}{\mathbb{R}}
\newcommand{\C}{\mathbb{C}}
\newcommand{\F}{\mathbb{F}}
\newcommand{\A}{\mathbb{A}}
\newcommand{\PP}{\mathbb{P}}
\newcommand{\E}{\mathbb{E}}
\newcommand{\eps}{\varepsilon}
\newcommand{\dd}{\,\mathrm{d}}
\begin{document}
\section*{UnsolvedMath ID: 20000234}
\textbf{Source problem code:} catalogue entry
\subsection{Definitions and mathematical statement}
The reconstructed, unverified Problem 9.2 asks for the sign of the plane Welschinger invariant with $d_1=2$ or $3$ real point conditions, predicting $(-1)^{(d-1)(d-2)/2+1}$ for sufficiently large degree. The extracted comparison before $0$ is corrupted; an eventual sign interpretation is supplied by the aid and companion context, not verified exact wording.

Companion Problem 9 specifies real irreducible rational degree-$d$ curves in $\PP^2$ through $d_1$ general real points and $d_2$ conjugate nonreal pairs, with $d_1+2d_2=3d-1$, counted with sign $(-1)^{s(C)}$ for the number $s(C)$ of isolated real nodes. Write this count $W_2(d,d_2)$ and set
\[
\delta_d=\frac{(d-1)(d-2)}2,\qquad
T_d=W_2\!\left(d,\left\lfloor\frac{3d-3}2\right\rfloor\right),\quad d\ge2.
\]
Then $d_1=2$ for odd $d$ and $d_1=3$ for even $d$. The intended eventual conjecture is
\[
\exists D\ \forall d\ge D:\quad (-1)^{\delta_d+1}T_d>0,
\]
not the assertion $T_d=(-1)^{\delta_d+1}$ and not necessarily an all-degree statement.

The attached unreviewed aid reports the termwise parity reduction
\[
(-1)^{\delta_d+1}T_d
=\#\{C:h(C)\text{ odd}\}-\#\{C:h(C)\text{ even}\},
\]
where $h(C)$ is the number of hyperbolic real nodes. If $\Gamma_C$ is the non-isolated real locus, it reports $\chi(\Gamma_C)=-h(C)$, so the normalized count is $-\sum_C(-1)^{\chi(\Gamma_C)}$. Thus a non-surjective injection from even-$h$ curves to odd-$h$ curves is a sufficient degeneration target, with this direction essential. It reports published exact values $T_d>0$ for $2\le d\le9$, implying failure of the all-degree strengthening at $d=2,5,6,9$ and the necessary threshold $D\ge10$ for an eventual statement.

Original Remark 9.3 asks for degeneration explanations of signs, potentially all curves having one sign or an injection between sign classes, and says signs on general real surfaces reflect topology. It lists floor-diagram formulas and open WDVV formulas as possible combinatorial approaches, noting that minimal $d_1$ or $d_2$ signs already have degeneration explanations. These are suggested approaches, not supplied proofs of the eventual conjecture.
\subsection{Short English statement}
Determine the eventual Welschinger sign when the number of real constraints is next to minimal. The aid converts it to a strict majority of odd-hyperbolic-node curves and reports low-degree exceptions forcing any eventual threshold to be at least ten.
\subsection{Sources}
\begin{itemize}
\item AIM \emph{Degenerations in algebraic geometry}, Problem 9.2 and Remark 9.3: \url{https://aimath.org/pastworkshops/degenalggeomproblems.pdf}. Canonical location: \texttt{aim-algebraic-geometry-notes.json}, notes[233]. The record is reconstructed/unverified, has no clean statement, and its verification is truncated. The remark's trailing ``10'' appears as a subsequent heading; no new problem is inferred from it.
\item Companion Problem 9, permanent ID 20000233, exact code \texttt{AIM-ALGEBRAIC\_GEOMETRY-0233}, canonical notes[232], attributes the surrounding sign question to ``Erwan'' only. Its verification corrects extracted $g_2$ to $d_2$ and identifies $\PP^2$. Remark 9.1 gives nonnegative sign for $d_2=0$ and sign $(-1)^{\delta_d}$ in the minimal-real-constraint case.
\item The present aid reports parity identities and published degree-$2$ through degree-$9$ values without identifying the publication. It is machine generated and not peer reviewed, with low novelty confidence; floor-diagram and open-WDVV references are also incomplete.
\item Local registry: \texttt{data/unsolvedmath/problems.json}, permanent catalogue entry, exact source code \texttt{catalogue entry}.
\end{itemize}
\end{document}
